%%\documentclass[a4paper]{article}
%%\usepackage{beamerarticle}
\documentclass[ignorenonframetext]{beamer}
\usepackage[backend=biber, autocite=plain, style=vancouver]{biblatex}
\addbibresource{referat.bib}
\usepackage[polish]{babel}
\usepackage[T1]{fontenc}
\usepackage[utf8]{inputenc}
\usepackage{amsmath}
\usepackage{amssymb}
\usepackage{csquotes}
\usefonttheme{professionalfonts}
\usepackage[scaled=0.92]{helvet}
\usepackage{newtxsf}
\usetheme{Warsaw}
\usecolortheme{seahorse}
\usecolortheme{rose}

\title{Jak zaprojektować metro?}
\subtitle{Matematyczne modele transportu publicznego i ich optymalizacja}
\author{Jan Pulkowski}
\institute{Uniwersytet Jagielloński}
\date{11.04.2026}

\DeclareMathOperator{\Ptn}{PTN}
\DeclareMathOperator{\Time}{time}
\DeclareMathOperator{\Inc}{inc}
\DeclareMathOperator{\Trans}{transfers}

\begin{document}

\frame{\titlepage}

\section{Wstęp}

\begin{frame}{Programowanie matematyczne}
  \begin{block}{Dane}
    \begin{itemize}
    \item Funkcja celu: \(f : A \to \mathbb{R}\)
    \item Funkcje ograniczeń: \(g_i : A \to \mathbb{R}\)
    \item Zbiór dopuszczalny:
      \[
        A_0 := \{x \in A : g_i(x) \ge 0 \text{ dla każdego } i\}
      \]
    \end{itemize}
  \end{block}
  \begin{block}{Cel}
    Znaleźć \(x_0 \in A_0\) takie, że
    \[
      f(x_0) = \max_{x \in A_0} f(x).
    \]
  \end{block}
\end{frame}
 O ile problemy programowania matematycznego można rozpatrywać na dowolnej dziedzinie \(A\), to znacznie prościej jest się ograniczyć do podzbiorów \(\mathbb{R}^{n}\). Znacznym uproszczeniem jest też przyjęcie dodatkowych założeń funkcjom celu i ograniczeń.
\begin{frame}{Programowanie liniowe i całkowitoliczbowe}
  \begin{block}{Programowanie liniowe}
    Szukamy optimum funkcji liniowej przy liniowych ograniczeniach:
    \[
      \max \{cx : Ax \le b,\ x \ge 0\}.
    \]
    gdzie \(x \in \mathbb{R}^n\), \(c \in \mathbb{R}^n\),
    \(A \in \mathbb{R}^{m \times n}\), \(b \in \mathbb{R}^m\).
  \end{block}
  \begin{block}{Programowanie całkowitoliczbowe}
    Dodatkowo wymagamy, aby wybrane lub wszystkie współrzędne wektora \(x\)
    były liczbami całkowitymi.
  \end{block}
\end{frame}

Jeżeli dodatkowo wymagamy, aby wszystkie (niektóre) współrzędne wektorów z \(X\) były liczbami całkowitymi, to będziemy mówić o (mieszanym) problemie programowania całkowitoliczbowego.

\begin{frame}{Jak to wykorzystać w komunikacji zbiorowej?}
  \begin{columns}[T]
    \begin{column}{0.48\textwidth}
      \begin{block}{Funkcja celu}
        \begin{itemize}
        \item koszty operatora,
        \item wygoda pasażerów,
        \item liczba podróży bezpośrednich.
        \end{itemize}
      \end{block}
    \end{column}
    \begin{column}{0.48\textwidth}
      \begin{block}{Ograniczenia}
        \begin{itemize}
        \item budżet,
        \item infrastruktura,
        \item minimalna jakość obsługi.
        \end{itemize}
      \end{block}
    \end{column}
  \end{columns}
  \vspace{0.3cm}
  Zajmiemy się przede wszystkim wygodą pasażerów przy ograniczonym budżecie i przepustowości sieci.
\end{frame}

\section{Model Schöbel--Scholl}

\begin{frame}{Model Schöbel--Scholl: podstawowe obiekty}
  \begin{itemize}
  \item \(\Ptn\) -- graf nieskierowany opisujący sieć transportową,
  \item \(S := V(\Ptn)\) -- zbiór przystanków,
  \item \(E := E(\Ptn)\) -- zbiór dopuszczalnych połączeń,
  \item \(\mathcal{L}\) -- zbiór możliwych linii,
  \item \(w_{st}\) -- zapotrzebowanie na podróże z \(s\) do \(t\),
  \item \(R := \{(s,t) \in S^2 : w_{st} > 0\}\) -- relacje z niezerowym popytem.
  \end{itemize}
\end{frame}

\begin{frame}{Model Schöbel--Scholl: co optymalizujemy?}
  Szukamy zbioru linii \(L \subseteq \mathcal{L}\) oraz ich częstotliwości \(f_l\), tak aby:
  \begin{itemize}
  \item każda relacja z \(R\) była obsłużona,
  \item łączny koszt nie przekraczał budżetu,
  \item podróż była możliwie wygodna dla pasażera.
  \end{itemize}
  \pause
  \begin{block}{Miara niewygody podróży}
    \[
      \Inc(P) = k_1\Time(P) + k_2\Trans(P)
    \]
    gdzie \(\Time(P)\) oznacza czas przejazdu, a \(\Trans(P)\) liczbę przesiadek.
  \end{block}
\end{frame}

\begin{frame}{Graf \emph{change \& go}}
  Aby modelować przesiadki, budujemy graf pomocniczy \(\mathcal{G} = (\mathcal{V}, \mathcal{E})\):
  \begin{itemize}
  \item \(\mathcal{V}_{CG}\): pary \emph{przystanek--linia},
  \item \(\mathcal{V}_{OD}\): początki i końce podróży,
  \item \(\mathcal{E}_{G}\): przejazdy po tej samej linii,
  \item \(\mathcal{E}_{C}\): przesiadki na tym samym przystanku,
  \item \(\mathcal{E}_{OD}\): wejście do sieci i wyjście z sieci.
  \end{itemize}
\end{frame}

Graf \emph{change \& go} formalnie:
\begin{flalign*}
  & \mathcal{G} = (\mathcal{V}, \mathcal{E}) \\
  & \mathcal{V} = \mathcal{V}_{CG} \cup \mathcal{V}_{OD} \\
  & \mathcal{V}_{CG} = \{ (s, l) \in S \times \mathcal{L} : s \in V(l)\} \\
  & \mathcal{V}_{OD} = \{ (s, \emptyset) : (s,t) \in R \lor (t,s) \in R\} \\
  & \mathcal{E} = \mathcal{E}_{C} \cup \mathcal{E}_{G} \cup \mathcal{E}_{OD} \\
  & \mathcal{E}_{C} = \{ ((s, l_1), (s, l_2)) \in \mathcal{V}_{CG} \times \mathcal{V}_{CG} \} \\
  & \mathcal{E}_{G} = \bigcup_{l \in \mathcal{L}} \{ ((s_1, l), (s_2, l)) \in \mathcal{V}_{CG} \times \mathcal{V}_{CG} : s_1s_2 \in E\} \\
  & \mathcal{E}_{OD} = \bigcup_{\substack{(s,t) \in R, \, l \in \mathcal{L} \\ s,t \in V(l)}} 
  \Bigl\{ \bigl( (s, \varnothing), (s, l) \bigr), \, \bigl( (t, l), (t, \varnothing) \bigr) \Bigr\} 
\end{flalign*}

\begin{frame}{Wagi krawędzi w grafie \emph{change \& go}}
  \[
    c_e =
    \begin{cases}
      0 & e \in \mathcal{E}_{OD}, \\
      k_1 t_{uv} & e = ((u,l),(v,l)) \in \mathcal{E}_{G}, \\
      k_2 & e \in \mathcal{E}_{C}.
    \end{cases}
  \]
  \begin{itemize}
  \item \(t_{uv}\) -- czas przejazdu między \(u\) i \(v\),
  \item \(k_1\) -- waga czasu,
  \item \(k_2\) -- koszt przesiadki.
  \end{itemize}
\end{frame}

Jeżeli interesuje nas tylko minimalizacja przesiadek, to możemy przyjąć, \(k_1 = 0, k_2 = 1\), jeżeli realistyczny czas przejazdu, to \(k_1 = 1\), a \(k_2\) to oszacowany czas przesiadki w tej samej jednostce, co czas pokonania krawędzi.

\begin{frame}{Ograniczenia w modelu Schöbel--Scholl}
  Model musi spełniać trzy grupy warunków:
  \begin{enumerate}
  \item \textbf{Obsługa popytu}: dla każdej relacji \((s,t) \in R\) istnieje poprawny przepływ.
  \item \textbf{Przepustowość}: liczba pasażerów nie przekracza pojemności i częstotliwości linii.
  \item \textbf{Budżet}: całkowity koszt eksploatacji mieści się w \(B\).
  \end{enumerate}
  Dodatkowo zmienne \(x_{st}^e\) i \(f_l\), reprezentujące przepływ pasażerów i częstotliwości linii, są całkowite.
\end{frame}

\begin{frame}{Model Schöbel--Scholl}
  \small
  \[
    \min \sum_{(s,t) \in R} \sum_{e \in \mathcal{E}} c_e x_{st}^e
  \]
  przy ograniczeniach:
  \begin{itemize}
  \item równania przepływu dla każdej relacji,
  \item ograniczenia pojemności,
  \item ograniczenia infrastruktury,
  \item ograniczenie budżetowe,
  \item całkowitość zmiennych.
  \end{itemize}
\end{frame}

Równania przepływu

Musimy znaleźć rozwiązania problemów sieci przepływowych
\[
  \forall (s, t) \in R \: \theta x_{st} = b_{st}
\]
\begin{itemize}
\item \(\theta \in \mathbb{Z}^{|\mathcal{V}| \times |\mathcal{E}|}\) -- macierz incydencji wierzchołków i krawędzi \(\mathcal{G}\)
\item \(x_{st} \in \mathbb{Z}^{|\mathcal{E}|}\) -- przepływ pasażerów z \(s\) do \(t\) po krawędzi \(e\)
\item \(b_{st} \in \mathbb{Z}^{|\mathcal{V}|}\), taki że:
  \[
    b_{st}^{i} =
    \begin{cases}
      w_{st} & i = (s, \emptyset)\\
      -w_{st} & i = (t, \emptyset)\\
      0 & \text{wpp.}
    \end{cases}
  \]
\end{itemize}

Ograniczenia Infranstruktury

\begin{enumerate}
\item \(\displaystyle \forall l \in \mathcal{L} \, \forall e \in \mathcal{E}_{l} \: x_{st}^{e} \le N \cdot f_{l} \)
\item \(\displaystyle \forall k \in E \: \sum_{l \in \mathcal{L} : k \in \mathcal{E}_{l}}f_{l} \le f_{k}^{max} \)
\end{enumerate}
\begin{itemize}
\item \(\mathcal{E}_{l} \subset \mathcal{E}_{G}\) -- krawędzie, po których prowadzi linia \(l\)
\item \(N\) -- pojemność jednego pojazdu
\item \(f_{k}^{max}\) -- maksymalna liczba pojazdów na danej krawędzi w analizowanym okresie
\end{itemize}

\begin{frame}{Wady modelu Schöbel--Scholl}
  \begin{alertblock}{Najważniejsze ograniczenia}
    \begin{itemize}
    \item Graf \emph{change \& go} może być bardzo duży.
    \item Model wymusza obsłużenie całego popytu.
    \item Nie rozróżnia podróży \enquote{akceptowalnych} i \enquote{nieakceptowalnych}.
    \end{itemize}
  \end{alertblock}
\end{frame}

Koszt ograniczony przez z góry ustalony budżet

\[
  \sum_{l \in \mathcal{L}} C_{l}f_{l} \le B
\]
\begin{itemize}
\item \(C_{l}\) -- koszt operacyjny 1 kursu na linii \(l\)
\item \(B\) -- ustalony budżet.
\end{itemize}

Pasażerowie jadą w 1 kawałku

\[
  \forall (s, t) \in R, e \in \mathcal{E}, l \in \mathcal{L} \: x_{st}^{e},f_{l} \in \mathbb{N}
\]

Pełen Model

\begin{alignat*}{3}
  \min \quad & \sum_{(s, t) \in R} \sum_{e \in \mathcal{E}} c_{e}x_{st}^{e} && \quad \text{takie że:} && \\
             & x_{st}^{e} \le N \cdot f_{l} && \quad && \forall l \in \mathcal{L}, \forall e \in \mathcal{E}_{l} \\
             & \theta x_{st} = b_{st}       && \quad && \forall (s, t) \in R \\
             & \sum_{l \in \mathcal{L}} C_{l}f_{l} \le B && \quad && \\
             & \sum_{\substack{l \in \mathcal{L} \\ k \in \mathcal{E}_{l}}} f_{l} \le f_{k}^{\max} && \quad && \forall k \in E \\
             & x_{st}^{e}, f_{l} \in \mathbb{N} && \quad && \forall (s, t) \in R, e \in \mathcal{E}, l \in \mathcal{L}
\end{alignat*}

\section{Model Nachtigall--Jerosch}

\begin{frame}{Dlaczego drugi model?}
  Chcemy modelu, który:
  \begin{itemize}
  \item lepiej oddaje zachowanie pasażerów,
  \item dopuszcza utratę części popytu,
  \item jest lżejszy obliczeniowo.
  \end{itemize}
\end{frame}

\begin{frame}{Model Nachtigall--Jerosch: główna idea}
  \begin{block}{Przewagi nad poprzednim modelem}
    \begin{itemize}
    \item Uwzględnia, że zbyt niewygodna podróż może zostać odrzucona przez pasażera.
    \item Ogranicza koszt pamięciowy i obliczeniowy reprezentacji sieci.
    \item Nie wymaga obsłużenia każdego pasażera.
    \end{itemize}
  \end{block}
  \pause
  \begin{exampleblock}{Zmiana celu}
    Zamiast minimalizować niewygodę każdej możliwej podróży, mierzymy jej
    jakość i wymagamy, by była wystarczająco wysoka.
  \end{exampleblock}
\end{frame}

\begin{frame}{Reprezentacja podróży}
  \begin{itemize}
  \item \(p^{st}_{(ab),l}\) -- fragment drogi z \(s\) do \(t\) realizowany odcinkiem \((ab)\) linii~\(l\),
  \item droga z \(s\) do \(t\) to konkatenacja takich fragmentów,
  \item \(x^{st}_{(ab),l}\) -- liczba pasażerów korzystających z danego fragmentu.
  \end{itemize}
\end{frame}

\begin{frame}{Jakość drogi}
  \begin{block}{Definicja}
    \[
      \omega(p^{st}) = \rho t_{st}^* - \Time(p^{st}) - \beta\Trans(p^{st}).
    \]
  \end{block}
  \begin{itemize}
  \item \(t_{st}^*\) -- najkrótszy możliwy czas przejazdu,
  \item \(\rho\) -- współczynnik tolerancji wydłużenia,
  \item \(\beta\) -- kara za przesiadkę.
  \end{itemize}
  \vspace{0.2cm}
  Im większa wartość \(\omega\), tym atrakcyjniejsza podróż. Podróże o~ujemnej wartości \(\omega\) uznajemy za nieakceptowalne.
\end{frame}

Wagi krawędzi

\[
  \omega^{st}_{(ab), l} =
  \begin{cases}
    -\Time_{(ab), l} & a = s, b \ne t  \\
    -\Time_{(ab), l} - \beta & a \ne s, b \ne t \\
    \rho t_{st}^{*}- \Time_{(ab), l} - \beta & a \ne s, b = t \\
    \rho t_{st}^{*} - \Time_{(ab), l} & a = s, b = t
  \end{cases}
\]

\begin{frame}{Dwa poziomy optymalizacji w modelu Nachtigall--Jerosch}
  \begin{enumerate}
  \item \textbf{Traffic assignment}: rozdział pasażerów na akceptowalne ścieżki.
  \item \textbf{Line planning}: dobór częstotliwości linii przy ograniczeniach kosztu i infrastruktury.
  \end{enumerate}
\end{frame}


Traffic Assignment

\begin{alignat*}{3}
  \max \: & \sum_{(s, t) \in R} \sum_{l \in \mathcal{L}} \sum_{a, b \in V(l)} \omega^{st}_{(ab), l} \varphi^{st}_{(ab), l} & \text{takie że:} \\
          & \sum_{l \in \mathcal{L}} \sum_{a \in V(l)} \varphi^{st}_{(at), l} \le w_{st} & \forall (s, t) \in R\\
          & \sum_{l \in \mathcal{L}} \sum_{a \in V(l)} \varphi^{st}_{(ax), l} - \sum_{l \in \mathcal{L}} \sum_{b \in V(l)} \varphi^{st}_{(xb), l} = 0 & \quad \forall (s, t) \in R, x \not \in \{s, t\} \\
          & \varphi^{st}_{(ab), l} \in \mathbb{Q}_{+}
\end{alignat*}


Line Planning

\begin{alignat*}{2}
  & \sum_{(s, t) \in R} \sum_{\substack{a, b \in V(l) \\ e \in l_{|ab}}} \varphi^{st}_{(xb), l} \le Nf_{l} & \quad \forall e \in E, l \in \mathcal{L} \\
  & \sum_{l \in \mathcal{L}} C_{l}f_{l} \le B \\
  & \sum_{\substack{l \in \mathcal{L} \\ e \in l}} f_{l} \le f^{max}(e) & \forall e \in E \\
  & \sum_{\substack{v \in \mathcal{V} \\ l \text{ serves } v}} f_{l} \le f^{max}(v) & \forall v \in V \\
  & f_{l} \in \mathbb{N} 
\end{alignat*}

\begin{frame}{Porównanie modeli}
  \small
  \begin{tabular}{p{0.33\textwidth} p{0.28\textwidth} p{0.28\textwidth}}
    \textbf{Cecha} & \textbf{Schöbel--Scholl} & \textbf{Nachtigall--Jerosch} \\
    \hline
    Reprezentacja podróży & graf \emph{change \& go} & uproszczony graf \\
    Obsługa popytu & pełna & częściowa \\
    Złożoność obliczeń & duża & mniejsza (wciąż NP) \\
  \end{tabular}
\end{frame}

Modele line planning pozwalają formalnie opisać kompromis między kosztem a wygodą pasażera. Model Schöbel--Scholl jest bardzo naturalny, ale kosztowny obliczeniowo. Model Nachtigall--Jerosch lepiej uwzględnia akceptowalność podróży i jest bardziej praktyczny. W rzeczywistym planowaniu trzeba łączyć jakość modelu z możliwością jego rozwiązania.

\begin{frame}{Dziękuję za uwagę}
  \nocite{*}
  \printbibliography
\end{frame}

\end{document}
